\section{Logical Architecture and Dependency Discipline}
\label{sec:hierarchy}

The Infrastructure tree is organized as a directed acyclic graph rather than a
flat collection of utilities.  Three independent roots have distinct roles:
\decl{SecurityParameter} defines \decl{Crypto.SecPar};
\decl{Probability} contains cost-free probability constructions; and
\decl{Cost.Model} defines the exact resource interface.  In particular,
computation modules no longer import the asymptotic layer merely to name a
security parameter.

\begin{figure}[t]
  \centering
  \resizebox{\linewidth}{!}{%
  \begin{tikzpicture}[node distance=7mm and 10mm]
    \node[rootbox] (sp) {Security\\Parameter};
    \node[rootbox, right=of sp] (prob) {Probability};
    \node[rootbox, right=of prob] (model) {Cost.Model};

    \node[archbox, below=of sp] (asy) {Asymptotic};
    \node[archbox, below=of model] (cost) {Cost writer, randomized\\semantics, path bounds};
    \node[archbox, below=of cost] (alg) {Typed algebra};
    \node[archbox, below left=of alg] (program) {Program};
    \node[archbox, below right=of alg] (oracle) {Oracle};
    \node[certbox, below=13mm of alg] (machine) {Complexity certificates\\and machines};
    \node[archbox, below=of machine] (game) {GameBased};
    \node[archbox, below=of game] (uc) {Typed UC execution};

    \draw[archarrow] (sp) -- (asy);
    \draw[archarrow] (model) -- (cost);
    \draw[archarrow] (cost) -- (alg);
    \draw[archarrow] (alg) -- (program);
    \draw[archarrow] (alg) -- (oracle);
    \draw[archarrow] (program) -- (machine);
    \draw[archarrow] (oracle) -- (machine);
    \draw[archarrow] (asy) -- (machine);
    \draw[archarrow] (machine) -- (game);
    \draw[archarrow] (game) -- (uc);
    \draw[archarrow] (machine.east) .. controls +(18mm,0) and +(18mm,0) .. (uc.east);
    \draw[archarrow] (sp.west) .. controls +(-14mm,-8mm) and +(-14mm,8mm) .. (uc.west);
  \end{tikzpicture}}
  \caption{Permitted logical flow.  Arrows denote admissible dependency
    directions, not necessarily direct Lean imports.  Probability,
    security parameters, and the base cost model are independent roots.}
  \label{fig:dependency-dag}
\end{figure}

The executable checker
\decl{scripts/check_infrastructure_imports.py} parses Infrastructure imports,
rejects cycles, and rejects edges that point upward in the declared subsystem
order.  Aggregation modules may import completed members of their own layer;
implementation modules must follow the finer order encoded by the checker.
This is a CI-enforced engineering invariant.  It is not presented as a
metatheorem about Lean's logic.

\begin{table}[t]
  \centering
  \caption{Authoritative objects by layer.  A higher row may export evidence
    consumed below it, but does not replace the semantic object it certifies.}
  \label{tab:layer-objects}
  \small
  \begin{tabular}{@{}>{\raggedright\arraybackslash}p{.20\linewidth}
      >{\raggedright\arraybackslash}p{.38\linewidth}
      >{\raggedright\arraybackslash}p{.34\linewidth}@{}}
    \toprule
    Layer & Authoritative object & Evidence or public view \\
    \midrule
    Cost & joint value--cost computation & support-wise path bound; cost erasure \\
    Algebra & typed exact operation handler & erased operation law; operation bound \\
    Program & one typed program and interpreter & execution relation; structural bound \\
    Oracle & one exact recursive interpreter & trace/count projections; composition bound \\
    Complexity & one dependent exact machine run & exact, measured, and polynomial certificates \\
    GameBased & cost-erased experiment & advantage, hardness, indistinguishability \\
    UC & typed closed-world kernel run & fuel/cost certificates; real/ideal decision distributions \\
    \bottomrule
  \end{tabular}
\end{table}

\subsection{Two orthogonal refinement directions}

The dependency DAG contains two conceptually different forms of refinement.
The semantic direction adds structure to exact execution:

\[
  \begin{gathered}
    \text{resource writer}\longrightarrow\text{typed operation handler}
    \\
    \longrightarrow\text{program/oracle/kernel execution}.
  \end{gathered}
\]

The certificate direction proves progressively more uniform statements without
changing that execution:

\[
  \begin{aligned}
    \text{path bound}&\longrightarrow\text{input-dependent exact budget}
    \\
    &\longrightarrow\text{uniform measured runtime}
      \longrightarrow\text{polynomial runtime}.
  \end{aligned}
\]

This distinction explains why asymptotics joins only at the complexity layer.
The definitions \decl{IsPolyBounded} and \decl{IsNegligible} do not participate
in exact evaluation.  Negligibility is stated using the eventual bound on
\(\lvert f(n)\rvert\), preventing a negative function from becoming negligible
merely because it lies below a positive threshold.

\subsection{Stable semantic boundaries}

Cost parameters do not enter ordinary games, advantages, or correctness
statements.  A security predicate instead selects an adversary cost model and
natural-number measure and then quantifies inhabitants of the resulting
PPT-machine interface whose exact run and claimed runtime also carry an
independent operational-admission witness.  Polynomial path annotations do not
construct that witness.  This preserves the distinction between the cost model
used to describe a construction and the independently selected operational
model used to admit adversaries.  Similarly, real and ideal UC worlds share
their environment, routing policy, corruption policy, and fuel convention;
only the installed protocol/adversary versus functionality/simulator components
differ.

The same pattern appears repeatedly: packaging may change as exact semantics
is exposed, but cost-erased probability statements remain the public boundary.
